% prompt-based-library-search.tex
% Prompt-Based Music Library Search Through Thermodynamic
% Ensemble Inference and Spectral Interference Matching
% Author: Kundai Farai Sachikonye
% Date: 2026

\documentclass[11pt,twocolumn,a4paper]{article}

% ─── Packages ────────────────────────────────────────────────────────
\usepackage[utf8]{inputenc}
\usepackage[T1]{fontenc}
\usepackage{lmodern}
\usepackage{amsmath,amssymb,amsthm,mathtools}
\usepackage{physics}
\usepackage{bm}
\usepackage{graphicx}
\usepackage{xcolor}
\usepackage{booktabs}
\usepackage{array}
\usepackage{multirow}
\usepackage{enumitem}
\usepackage{float}
\usepackage{caption}
\usepackage{subcaption}
\usepackage[colorlinks=true,linkcolor=blue!60!black,citecolor=blue!60!black,urlcolor=blue!60!black]{hyperref}
\usepackage[margin=1.8cm,top=2.2cm,bottom=2.2cm]{geometry}
\usepackage{natbib}
\usepackage{fancyhdr}
\usepackage{titlesec}
\usepackage{microtype}
\usepackage{algorithm}
\usepackage{algpseudocode}
\usepackage{listings}
\usepackage{tabularx}

% ─── Listings Style ─────────────────────────────────────────────────
\lstdefinestyle{glsl}{
  language=C,
  basicstyle=\ttfamily\scriptsize,
  keywordstyle=\color{blue!70!black}\bfseries,
  commentstyle=\color{green!50!black}\itshape,
  stringstyle=\color{red!60!black},
  numbers=left,
  numberstyle=\tiny\color{gray},
  numbersep=4pt,
  frame=single,
  breaklines=true,
  columns=fullflexible,
  morekeywords={vec2,vec3,vec4,mat2,mat3,mat4,sampler2D,sampler3D,
    uniform,varying,in,out,layout,location,precision,
    highp,mediump,lowp,texture,gl_FragColor,void,float,int,
    ivec3,uvec3,bool}
}

% ─── Theorem Environments ────────────────────────────────────────────
\newtheorem{theorem}{Theorem}[section]
\newtheorem{lemma}[theorem]{Lemma}
\newtheorem{corollary}[theorem]{Corollary}
\newtheorem{proposition}[theorem]{Proposition}
\newtheorem{definition}[theorem]{Definition}
\newtheorem{remark}[theorem]{Remark}
\newtheorem{axiom}{Axiom}
\newtheorem{example}[theorem]{Example}

% ─── Custom Commands ─────────────────────────────────────────────────
\newcommand{\kB}{k_{\mathrm{B}}}
\newcommand{\Ocat}{\hat{O}_{\mathrm{cat}}}
\newcommand{\Ophys}{\hat{O}_{\mathrm{phys}}}
\newcommand{\Sk}{S_k}
\newcommand{\St}{S_t}
\newcommand{\Se}{S_e}
\newcommand{\R}{\mathbb{R}}
\newcommand{\N}{\mathbb{N}}
\newcommand{\Z}{\mathbb{Z}}
\newcommand{\CC}{\mathbb{C}}
\newcommand{\Lib}{\mathcal{L}}
\newcommand{\Ens}{\mathcal{E}}
\newcommand{\Part}{\mathcal{P}}
\newcommand{\Cat}{\mathcal{C}}
\newcommand{\Vis}{\mathcal{V}}
\newcommand{\Query}{\mathcal{Q}}
\newcommand{\Spec}{\bm{\Phi}}
\DeclareMathOperator{\Tr}{Tr}
\DeclareMathOperator{\diag}{diag}
\DeclareMathOperator{\sgn}{sgn}

% ─── Page Style ──────────────────────────────────────────────────────
\pagestyle{fancy}
\fancyhf{}
\fancyhead[L]{\small\textit{Prompt-Based Library Search via Thermodynamic Inference}}
\fancyhead[R]{\small\thepage}
\renewcommand{\headrulewidth}{0.4pt}

% ─── Title Section Formatting ────────────────────────────────────────
\titleformat{\section}{\large\bfseries}{\thesection.}{0.5em}{}
\titleformat{\subsection}{\normalsize\bfseries}{\thesubsection.}{0.5em}{}
\titleformat{\subsubsection}{\normalsize\itshape}{\thesubsubsection.}{0.5em}{}

% ═════════════════════════════════════════════════════════════════════
\begin{document}

\twocolumn[
\begin{@twocolumnfalse}
\begin{center}

{\LARGE\bfseries Prompt-Based Music Library Search Through\\[4pt]
Thermodynamic Ensemble Inference and\\[4pt]
Spectral Interference Matching}

\vspace{12pt}

{\large Kundai Farai Sachikonye}

\vspace{6pt}

{\small\itshape Independent Researcher \\[0.3em]
\texttt{kundai.sachikonye@bitspark.com}}

\vspace{6pt}

{\small April 2026}

\vspace{14pt}

\begin{minipage}{0.92\textwidth}
\textbf{Abstract.}
We present a music search and recommendation framework in which the
user's listening library is treated as a thermodynamic gas ensemble, each
track is a molecule characterised by its categorical audio state, and
all search, matching, and recommendation operations reduce to
well-posed problems in statistical mechanics. The framework rests on
three results from categorical audio theory
\citep{sachikonye2026categorical}: (i)~audio signals in bounded phase
space are equivalent to oscillatory systems with partition coordinates
$(n, \ell, m, s)$ and S-entropy coordinates $(\Sk, \St, \Se)$;
(ii)~oscillatory systems have spectra; (iii)~spectra are rendered by
GPU fragment shaders, so rendering is computation (the Triple
Equivalence). Because every track is an oscillator ensemble expressible
as a spectrum, the similarity between two tracks is their
\emph{interference visibility}---the magnitude of the coherent
superposition of their phase spectra---not a distance in an embedding
space. The listening library, as a gas ensemble, possesses
well-defined thermodynamic quantities: temperature (taste breadth),
partition function (statistical description of preference), Boltzmann
distribution (probability of enjoying a candidate track), free energy
(equilibrium preference state), and entropy (predictability of taste).
Finding the next track is sampling from the Boltzmann distribution;
discovering new music is exploring the tails; refining taste is
annealing the ensemble temperature. A natural-language interface is
provided by a domain-specific language model trained via structured
knowledge distillation (the Purpose framework) on the user's
categorical observation library, translating between human prompts and
thermodynamic operations on the ensemble. The GPU fragment shader that
renders the liquid-distortion visualisation of album artwork
simultaneously computes the Boltzmann weights, interference
visibilities, and partition function of the ensemble, because these are
the same spectral computation expressed in three equivalent
representations. We derive the complete mathematical framework from
first principles, prove that interference matching recovers all
information available to embedding-based methods while additionally
capturing cross-genre structural similarity inaccessible to learned
representations, and specify the five-component architecture: (1)~categorical
observation pipeline, (2)~thermodynamic ensemble engine,
(3)~interference matching shader, (4)~Purpose-trained language model
interface, and (5)~streaming API integration. The system requires no
training data, has no cold-start problem, provides per-channel
interpretability of every similarity judgment, and improves with use
through ensemble enrichment rather than collaborative filtering.

\vspace{8pt}
\noindent\textbf{Keywords:} thermodynamic ensemble, interference
matching, categorical audio, S-entropy, partition function, Boltzmann
distribution, music recommendation, spectral similarity, GPU
computation, knowledge distillation, prompt-based search
\end{minipage}

\vspace{18pt}
\end{center}
\end{@twocolumnfalse}
]


% ═════════════════════════════════════════════════════════════════════
%                    PART I: FOUNDATIONS
% ═════════════════════════════════════════════════════════════════════

\section{Introduction}
\label{sec:introduction}

Music recommendation systems face a fundamental representation problem.
Collaborative filtering methods \citep{koren2009matrix,
covington2016youtube} represent tracks as rows in a user-interaction
matrix and compute similarity through co-occurrence statistics. The
track itself---its physical audio content---is absent from the
representation. Content-based methods extract feature vectors (MFCCs
\citep{logan2000mel}, learned embeddings \citep{wu2023clap,
radford2023whisper}) and compute distances in feature space. The
features are proxies for perceptual similarity, learned from data,
entangled across dimensions, and opaque to interpretation.

Both approaches share a structural limitation: they represent each
track as a \emph{point} in a high-dimensional space and compute
similarity as \emph{distance} between points. This reduces
recommendation to a nearest-neighbour search---a geometric problem. We
demonstrate that this geometric framing discards information that is
physically present in the audio signal.

We propose an alternative: represent each track as an \emph{oscillator
ensemble} in categorical phase space, and compute similarity as
\emph{interference} between ensembles. This is not a metaphor. The
Categorical Audio Transport (CAT) framework
\citep{sachikonye2026categorical} proves that audio signals in bounded
phase space are equivalent to collections of oscillatory modes, each
addressed by partition coordinates $(n, \ell, m, s)$ and projecting
onto S-entropy coordinates $(\Sk, \St, \Se)$. Oscillators have
spectra. Spectra interfere. The interference visibility between two
tracks' spectra is a physically grounded measure of their structural
similarity.

The user's listening library, as a collection of such oscillator
ensembles, constitutes a thermodynamic gas ensemble with well-defined
statistical mechanical properties. This recasts music recommendation as
a sampling problem in statistical mechanics rather than a search problem
in metric space.

\subsection{Contributions}

\begin{enumerate}[leftmargin=*,itemsep=2pt]
\item We define the \emph{listening library as a thermodynamic
  ensemble} (Section~\ref{sec:ensemble}) with derivations of
  temperature, partition function, Boltzmann distribution, free energy,
  and entropy from the categorical audio states of observed tracks.

\item We prove that \emph{interference visibility} between categorical
  phase spectra (Section~\ref{sec:interference}) recovers all pairwise
  similarity information available to inner-product methods while
  additionally capturing structural relationships across genre
  boundaries that embedding-based methods cannot represent.

\item We establish the \emph{Rendering--Computing--Observing Identity}
  (Section~\ref{sec:rco}): the GPU fragment shader that renders visual
  effects on album artwork simultaneously evaluates the Boltzmann
  weights and interference visibilities of the ensemble, because these
  are the same spectral computation in three equivalent representations.

\item We specify the \emph{Purpose-trained language model interface}
  (Section~\ref{sec:purpose}): a domain-specific model trained via
  structured knowledge distillation on the user's categorical
  observation library, translating natural-language prompts into
  thermodynamic operations on the ensemble.

\item We derive the \emph{ensemble enrichment dynamics}
  (Section~\ref{sec:enrichment}): each new track observation refines
  the thermodynamic description of the library, improving all
  subsequent recommendations without retraining.
\end{enumerate}

\subsection{Relation to Prior Work}

The framework builds on three prior results:

\textbf{Categorical Audio Transport} \citep{sachikonye2026categorical}
establishes the partition coordinate system $(n, \ell, m, s)$, the
S-entropy coordinates $(\Sk, \St, \Se)$, the Triple Equivalence
between oscillatory, categorical, and partitional descriptions, and the
commutation relation $[\Ocat, \Ophys] = 0$ between categorical and
physical observables.

\textbf{Ray-Tracing as Cellular Computation}
\citep{sachikonye2026raytracing} proves the Triple Observation
Identity: a single render pass simultaneously computes optical,
chromatographic, and circuit observations from one partition state read.
This establishes that rendering IS computation IS observation for any
system expressible as partition coordinates.

\textbf{Purpose Framework} provides the structured knowledge
distillation pipeline for training domain-specific language models from
extracted observations, using stratified query generation across
knowledge depth and cognitive levels \citep{bloom1956taxonomy,
hinton2015distilling, hu2021lora}.


% ═════════════════════════════════════════════════════════════════════
\section{Categorical Audio State}
\label{sec:categorical}

We summarise the results from \citet{sachikonye2026categorical} required
for the present work.

\begin{definition}[Track as Oscillator Ensemble]
\label{def:track-ensemble}
An audio track $\mathbf{x}(t)$ of duration $T$ sampled at rate $f_s$,
analysed with $M$ spectral modes, is represented as an ensemble of $M$
oscillators, each characterised by:
\begin{enumerate}[label=(\roman*),leftmargin=*]
\item Partition coordinates $(n_k, \ell_k, m_k, s_k)$ for $k = 1,
  \ldots, M$, encoding the categorical state of mode $k$.
\item S-entropy coordinates $(\Sk, \St, \Se)$ computed from the
  aggregate signal as:
  \begin{align}
    \Sk &= \kB \ln\!\left(\frac{|\delta\phi| + \phi_0}{\phi_0}\right)
    \label{eq:sk} \\
    \St &= \kB \ln\!\left(\frac{\tau}{\tau_0}\right)
    \label{eq:st} \\
    \Se &= \kB \ln\!\left(\frac{E + E_0}{E_0}\right)
    \label{eq:se}
  \end{align}
  where $\delta\phi$ is the phase deviation from a reference
  oscillator, $\tau$ is the categorical period, $E$ is the
  instantaneous signal energy, and $\phi_0$, $\tau_0$, $E_0$ are
  reference values.
\item A phase spectrum $\Spec = (\Phi_1, \ldots, \Phi_N)$ where
  $\Phi_k$ is the accumulated phase contribution of oscillator class
  $k$ over the track duration:
  \begin{equation}
    \Phi_k = \int_0^T \omega_k(t)\, \tau_{\mathrm{p}}^{(k)}(t)\, dt
    \label{eq:phase-accumulation}
  \end{equation}
  where $\omega_k(t)$ is the instantaneous frequency and
  $\tau_{\mathrm{p}}^{(k)}(t)$ is the interaction time of oscillator
  class $k$ at time $t$.
\end{enumerate}
\end{definition}

\begin{remark}
The phase spectrum is the track's categorical fingerprint. Two tracks
with identical phase spectra are categorically indistinguishable
regardless of their physical waveforms. The phase spectrum encodes the
cumulative interaction history of every oscillator class, analogous to
the accumulated phase shift along a ray path in the cellular
computation framework \citep{sachikonye2026raytracing}.
\end{remark}

\begin{definition}[Track Observation]
\label{def:observation}
A \emph{track observation} is the structured tuple:
\begin{equation}
  \mathcal{O}_i = \bigl(\bar{\mathbf{S}}_i,\, \sigma_{\mathbf{S},i},\,
  \mathbf{h}_i,\, \Spec_i,\, n_i^*\bigr)
  \label{eq:observation}
\end{equation}
where $\bar{\mathbf{S}}_i = (\langle \Sk \rangle, \langle \St \rangle,
\langle \Se \rangle)$ is the mean S-entropy over the track duration,
$\sigma_{\mathbf{S},i}$ is the standard deviation of the S-entropy
trajectory, $\mathbf{h}_i$ is the normalised partition depth histogram,
$\Spec_i$ is the phase spectrum, and $n_i^*$ is the dominant partition
depth.
\end{definition}

\begin{definition}[Segmented Observation]
\label{def:segmented}
A track of duration $T$ is partitioned into $L$ segments of duration
$\Delta T = T / L$. Each segment $\ell \in \{1, \ldots, L\}$ yields an
independent observation $\mathcal{O}_i^{(\ell)}$. The segmented
observation of the track is:
\begin{equation}
  \tilde{\mathcal{O}}_i =
  \bigl\{\mathcal{O}_i^{(1)}, \ldots, \mathcal{O}_i^{(L)}\bigr\}
  \label{eq:segmented}
\end{equation}
\end{definition}

\begin{remark}
Segmented observation enables temporal matching: the intro of track $A$
may interfere constructively with the breakdown of track $B$. This is
inaccessible to whole-track representations.
\end{remark}


% ═════════════════════════════════════════════════════════════════════
\section{The Listening Library as a Thermodynamic Ensemble}
\label{sec:ensemble}

\begin{definition}[Listening Library]
\label{def:library}
A user's listening library is the ordered collection of track
observations accumulated over time:
\begin{equation}
  \Lib = \bigl\{\mathcal{O}_1, \mathcal{O}_2, \ldots,
  \mathcal{O}_M\bigr\}
  \label{eq:library}
\end{equation}
where $M$ is the number of distinct tracks observed.
\end{definition}

\begin{proposition}[Library as Gas Ensemble]
\label{prop:gas-ensemble}
The listening library $\Lib$ satisfies the axioms of a classical
thermodynamic gas ensemble:
\begin{enumerate}[label=(\roman*),leftmargin=*]
\item Each track $\mathcal{O}_i$ is a molecule with internal degrees
  of freedom (partition coordinates) and translational energy
  ($\Se_i$).
\item The ensemble occupies a bounded region of S-entropy space
  $\mathcal{S} \subset \R_{\geq 0}^3$ (since all S-entropy
  coordinates are non-negative and bounded above by the dynamic range
  of the audio system).
\item The molecules interact through interference (Section~\ref{sec:interference}), with interaction strength
  determined by phase spectrum overlap.
\end{enumerate}
\end{proposition}

\begin{proof}
\textit{Step 1: Molecular identification.}
By Definition~\ref{def:track-ensemble}, each track is an ensemble of
$M_k$ oscillators with partition coordinates and accumulated phases.
The track observation $\mathcal{O}_i$ is therefore a point in
$\mathcal{S} \times \R^N$ (S-entropy space cross phase space),
analogous to a molecule with position $\bar{\mathbf{S}}_i$ and
internal state $\Spec_i$.

\textit{Step 2: Boundedness.}
The S-entropy coordinates are computed via logarithms of positive
ratios (Eqs.~\eqref{eq:sk}--\eqref{eq:se}). The audio system has
finite dynamic range (bit depth), finite bandwidth (sampling rate), and
finite temporal extent. Therefore $\Sk$, $\St$, $\Se$ are each bounded
above, and the ensemble occupies a bounded region of $\R^3$.

\textit{Step 3: Interactions.}
The interference visibility between tracks $i$ and $j$
(Definition~\ref{def:visibility}) provides a pairwise interaction
potential, establishing the ensemble as an interacting gas rather than
an ideal gas.
\end{proof}

We now derive the thermodynamic quantities of the library ensemble.

\begin{definition}[Ensemble Temperature]
\label{def:temperature}
The temperature of the library ensemble is defined by the equipartition
theorem applied to the S-entropy kinetic energy:
\begin{equation}
  T_{\Lib} = \frac{2}{3\kB} \langle E \rangle
  = \frac{2}{3\kB} \cdot \frac{1}{M} \sum_{i=1}^{M} \Se_i
  \label{eq:temperature}
\end{equation}
where $\Se_i = \langle \Se \rangle_i$ is the mean energetic entropy of
track $i$.
\end{definition}

\begin{remark}
High ensemble temperature corresponds to a listener with broad
energetic range across their library (eclectic taste). Low temperature
corresponds to a library concentrated at specific energy levels (narrow
preference).
\end{remark}

\begin{definition}[Partition Function]
\label{def:partition-function}
The canonical partition function of the library ensemble at inverse
temperature $\beta = 1/(\kB T_{\Lib})$ is:
\begin{equation}
  Z(\beta) = \sum_{i=1}^{M}
  \exp\!\left(-\beta\, \|\bar{\mathbf{S}}_i\|^2\right)
  \label{eq:partition-function}
\end{equation}
where $\|\bar{\mathbf{S}}_i\|^2 = \Sk_i^2 + \St_i^2 + \Se_i^2$ is
the squared magnitude of the mean S-entropy vector of track $i$.
\end{definition}

\begin{definition}[Boltzmann Distribution]
\label{def:boltzmann}
The probability that a candidate track with S-entropy $\mathbf{S}$ is
compatible with the library ensemble is:
\begin{equation}
  P(\mathbf{S}\,|\,\Lib) = \frac{1}{Z(\beta)}
  \exp\!\left(-\beta\, d_G(\mathbf{S},\,
  \bar{\mathbf{S}}_{\Lib})^2\right)
  \label{eq:boltzmann}
\end{equation}
where $d_G$ is the geodesic distance on the S-entropy manifold under
the groove metric:
\begin{equation}
  d_G(\mathbf{S}_a, \mathbf{S}_b)^2 = g_{kk}\,\Delta\Sk^2
  + 2g_{kt}\,\Delta\Sk\,\Delta\St + g_{tt}\,\Delta\St^2
  + g_{ee}\,\Delta\Se^2
  \label{eq:geodesic}
\end{equation}
with metric tensor components $g_{kk} = (\Sk + \epsilon)^{-1}$,
$g_{tt} = (\St + \epsilon)^{-1}$, $g_{ee} = (\Se + \epsilon)^{-1}$,
$g_{kt} = \tfrac{1}{2}\sqrt{\Sk\,\St}$, and
$\bar{\mathbf{S}}_{\Lib} = M^{-1}\sum_i \bar{\mathbf{S}}_i$ is the
ensemble centroid.
\end{definition}

\begin{definition}[Free Energy]
\label{def:free-energy}
The Helmholtz free energy of the library ensemble is:
\begin{equation}
  F = -\kB T_{\Lib} \ln Z(\beta)
  \label{eq:free-energy}
\end{equation}
The minimum of $F$ identifies the equilibrium preference state: the
S-entropy region the listener most naturally gravitates toward.
\end{definition}

\begin{definition}[Taste Entropy]
\label{def:taste-entropy}
The entropy of the library ensemble is:
\begin{equation}
  S_{\mathrm{taste}} = -\kB \sum_{i=1}^{M} P_i \ln P_i
  \label{eq:taste-entropy}
\end{equation}
where $P_i = Z^{-1} \exp(-\beta\, \|\bar{\mathbf{S}}_i\|^2)$.
High taste entropy indicates an unpredictable listener; low taste
entropy indicates a listener whose next preference is highly
constrained.
\end{definition}

\begin{theorem}[Recommendation as Boltzmann Sampling]
\label{thm:recommendation}
Finding the optimal next track for a listener with library $\Lib$ is
equivalent to sampling from the Boltzmann distribution
$P(\mathbf{S}\,|\,\Lib)$ at the ensemble temperature $T_{\Lib}$.
Specifically:
\begin{enumerate}[label=(\roman*),leftmargin=*]
\item ``More of the same'' corresponds to sampling near the mode of
  $P(\mathbf{S}\,|\,\Lib)$ (maximum likelihood).
\item ``Surprise me'' corresponds to sampling from the tails of
  $P(\mathbf{S}\,|\,\Lib)$ at elevated temperature $T' > T_{\Lib}$.
\item ``Something darker'' corresponds to sampling from
  $P(\mathbf{S}\,|\,\Lib)$ with a shifted constraint
  $\Se > \bar{\Se}_{\Lib} + \delta$.
\item ``Like Tuesday evening'' corresponds to sampling from the
  Boltzmann distribution of the temporal sub-ensemble
  $\Lib_{\mathrm{Tue}} \subset \Lib$.
\end{enumerate}
\end{theorem}

\begin{proof}
Each operation is a standard statistical mechanical manipulation of the
ensemble:

\textit{(i)} Maximum-likelihood sampling: the mode of the Boltzmann
distribution is the state minimising $d_G(\mathbf{S},
\bar{\mathbf{S}}_{\Lib})$, which is the ensemble centroid.

\textit{(ii)} Temperature elevation: increasing $\beta^{-1}$ broadens
the distribution, assigning higher probability to states far from the
centroid. This is simulated annealing in reverse---heating the ensemble
to explore.

\textit{(iii)} Constrained sampling: adding the constraint $\Se >
\bar{\Se}_{\Lib} + \delta$ restricts the Boltzmann distribution to a
half-space in S-entropy space, which remains a valid (unnormalised)
Gibbs measure.

\textit{(iv)} Sub-ensemble: the temporal sub-ensemble
$\Lib_{\mathrm{Tue}}$ defines its own partition function $Z_{\mathrm{Tue}}$,
temperature $T_{\mathrm{Tue}}$, and Boltzmann distribution. Sampling
from this sub-ensemble is identical to sampling from the full ensemble
conditioned on the temporal index.
\end{proof}


% ═════════════════════════════════════════════════════════════════════
\section{Interference Matching}
\label{sec:interference}

\subsection{Interference Visibility}

\begin{definition}[Interference Visibility]
\label{def:visibility}
The interference visibility between two tracks with phase spectra
$\Spec_A = (\Phi_{A,1}, \ldots, \Phi_{A,N})$ and
$\Spec_B = (\Phi_{B,1}, \ldots, \Phi_{B,N})$ is:
\begin{equation}
  V_{AB} = \left|\frac{1}{N} \sum_{k=1}^{N}
  e^{i(\Phi_{A,k} - \Phi_{B,k})}\right|
  \label{eq:visibility}
\end{equation}
where $N$ is the number of oscillator classes. $V_{AB} \in [0, 1]$
with $V_{AB} = 1$ indicating perfect constructive interference
(categorical identity) and $V_{AB} = 0$ indicating complete
decoherence (categorical independence).
\end{definition}

\begin{proposition}[Visibility as Complex Inner Product]
\label{prop:inner-product}
Define the complex unit phasor representation of a track:
$\mathbf{z}_A = (e^{i\Phi_{A,1}}, \ldots, e^{i\Phi_{A,N}}) \in
\CC^N$. Then:
\begin{equation}
  V_{AB} = \frac{1}{N}
  \bigl|\langle \mathbf{z}_A, \mathbf{z}_B \rangle\bigr|
  = \frac{1}{N}
  \left|\sum_{k=1}^{N} \overline{z_{A,k}}\, z_{B,k}\right|
  \label{eq:inner-product}
\end{equation}
This is the magnitude of the Hermitian inner product on $\CC^N$,
normalised by $N$.
\end{proposition}

\begin{proof}
Direct expansion: $\overline{z_{A,k}} \cdot z_{B,k} =
e^{-i\Phi_{A,k}} \cdot e^{i\Phi_{B,k}} = e^{i(\Phi_{B,k} -
\Phi_{A,k})}$. Summing and taking the magnitude yields
Eq.~\eqref{eq:visibility} with the sign convention reversed (which
does not affect the magnitude).
\end{proof}

\begin{remark}
Equation~\eqref{eq:inner-product} shows that interference visibility
is mathematically equivalent to cosine similarity in $\CC^N$. The
novelty is not in the comparison operation but in the
\emph{representation}: the phase spectrum $\Spec$ encodes
second-order structural properties (spectral complexity, temporal
granularity, energetic dynamics, inter-band coupling) that are absent
from standard audio feature vectors.
\end{remark}

\subsection{Per-Channel Decomposition}

\begin{definition}[Channel Coherence]
\label{def:channel-coherence}
The coherence of oscillator class $k$ between tracks $A$ and $B$ is:
\begin{equation}
  c_k = \cos(\Phi_{A,k} - \Phi_{B,k})
  \label{eq:channel-coherence}
\end{equation}
with $c_k = 1$ indicating phase alignment (constructive interference
in class $k$), $c_k = -1$ indicating anti-phase (destructive
interference), and $c_k = 0$ indicating quadrature (independence).
\end{definition}

\begin{theorem}[Interpretability of Interference]
\label{thm:interpretability}
The channel coherence vector
$\mathbf{c} = (c_1, \ldots, c_N)$ provides a complete,
per-channel decomposition of the similarity judgment. Specifically:
\begin{enumerate}[label=(\roman*),leftmargin=*]
\item The global visibility is recovered: $V_{AB}^2 = N^{-2}
  \bigl(\sum_k c_k\bigr)^2 + N^{-2}\bigl(\sum_k
  \sin(\Phi_{A,k} - \Phi_{B,k})\bigr)^2$.
\item Each channel $c_k$ is independently interpretable: it identifies
  whether oscillator class $k$ is aligned, opposed, or independent
  between the two tracks.
\item The decomposition is invertible: given $\mathbf{c}$ and
  $V_{AB}$, the phase difference vector $\Delta\Spec = \Spec_A -
  \Spec_B$ is determined up to a global phase.
\end{enumerate}
\end{theorem}

\begin{proof}
\textit{(i)} follows from $V_{AB} = N^{-1}|\sum_k e^{i\Delta\Phi_k}|
= N^{-1}\sqrt{(\sum_k \cos\Delta\Phi_k)^2 + (\sum_k
\sin\Delta\Phi_k)^2}$.

\textit{(ii)} follows from the independence of the oscillator classes:
each $\Phi_k$ is accumulated from a distinct frequency band, so
$c_k$ depends only on the phase relationship in class $k$.

\textit{(iii)} follows from $\Delta\Phi_k = \arccos(c_k)$ (with sign
ambiguity resolved by $\sin\Delta\Phi_k$).
\end{proof}

\begin{corollary}[Embedding Methods Lack Interpretability]
\label{cor:embedding-opaque}
Cosine similarity in a learned embedding space $\R^d$ yields a scalar
$s_{AB} = \langle \mathbf{e}_A, \mathbf{e}_B \rangle / (\|\mathbf{e}_A\|
\|\mathbf{e}_B\|)$ with no per-dimension physical interpretation
because the embedding dimensions are entangled by training. Interference
visibility in $\CC^N$ yields the same scalar \emph{plus} the
decomposition $\mathbf{c}$, because the oscillator classes are
physically independent.
\end{corollary}

\subsection{Cross-Genre Structural Similarity}

\begin{theorem}[Structural Similarity Across Genres]
\label{thm:cross-genre}
Let $A$ be a track from genre $\alpha$ and $B$ a track from genre
$\beta \neq \alpha$. If the S-entropy trajectories and partition depth
distributions of $A$ and $B$ are statistically similar, then:
\begin{enumerate}[label=(\roman*),leftmargin=*]
\item The interference visibility $V_{AB}$ is high ($V_{AB} > 0.7$).
\item The Boltzmann probability $P(\mathbf{S}_B\,|\,\Lib)$ is high
  when $A \in \Lib$.
\item Embedding-based similarity $s_{AB}$ may be low if the
  embedding was trained on genre-labelled data, because the training
  objective separates genres in embedding space.
\end{enumerate}
\end{theorem}

\begin{proof}
\textit{(i)} If the S-entropy trajectories of $A$ and $B$ are
statistically similar, their phase spectra $\Spec_A$ and $\Spec_B$ are
accumulated from similar instantaneous frequencies and interaction
times (Eq.~\eqref{eq:phase-accumulation}), yielding small phase
differences $\Delta\Phi_k$ and therefore high $V_{AB}$ by
Eq.~\eqref{eq:visibility}.

\textit{(ii)} Statistical similarity of S-entropy trajectories implies
$d_G(\bar{\mathbf{S}}_B, \bar{\mathbf{S}}_A)$ is small, and since $A
\in \Lib$, the Boltzmann distribution assigns high probability to the
region near $\bar{\mathbf{S}}_A$.

\textit{(iii)} Learned embeddings trained with genre labels as
supervision (or with data biased by genre co-occurrence) allocate
orthogonal subspaces to distinct genres. If $\alpha$ and $\beta$ are
sufficiently distinct in the training distribution, $\|\mathbf{e}_A -
\mathbf{e}_B\|$ is large even when the underlying categorical
structure is identical.
\end{proof}

\begin{example}
A jazz piano trio and a neurofunk drum \& bass track may share: high
$\Sk$ (complex harmonic structure), similar $\St$ fluctuation (swung
micro-timing), compatible partition depth evolution, and parallel groove
geodesics. Their interference visibility is high. Their CLAP embedding
distance is large because CLAP was trained to separate ``jazz'' from
``drum and bass'' in text-audio space. The categorical representation
captures the structural connection; the learned representation cannot.
\end{example}


% ═════════════════════════════════════════════════════════════════════
\section{Rendering--Computing--Observing Identity}
\label{sec:rco}

\begin{theorem}[Spectral Computation Identity]
\label{thm:rco}
Let $\mathcal{O}_i$ be a track observation with phase spectrum
$\Spec_i$. The GPU fragment shader that renders the visual distortion
of the album artwork simultaneously computes:
\begin{enumerate}[label=(\roman*),leftmargin=*]
\item The Boltzmann weight $w_i = \exp(-\beta\,
  \|\bar{\mathbf{S}}_i\|^2)$ from the S-entropy coordinates (which
  parameterise the shader uniforms).
\item The interference visibility $V_{ij}$ for any pair $(i, j)$ via
  superposition of the rendered wave fields.
\item The partition function contribution $Z_i = w_i$ as a partial
  sum.
\end{enumerate}
These are not three separate computations performed in sequence but a
single spectral evaluation expressed in three equivalent
representations.
\end{theorem}

\begin{proof}
\textit{Step 1: Shader as spectral evaluator.}
The fragment shader evaluates, at each pixel $(u, v)$, a sum of
sinusoidal displacements:
\begin{equation}
  D(u, v, t) = \sum_{k=1}^{N} A_k \sin\!\bigl(2\pi f_k\,
  \mathbf{n}_k \cdot (u, v) + \Phi_k + \omega_k t\bigr)
  \label{eq:shader-displacement}
\end{equation}
where $A_k$, $f_k$, $\omega_k$ are determined by the audio data
(bass, mid, treble, volume) and $\Phi_k$ is the accumulated phase of
oscillator class $k$. This is a spectral synthesis---the shader is
evaluating the spectrum of the track's oscillator ensemble.

\textit{Step 2: Boltzmann weight from shader uniforms.}
The S-entropy coordinates $(\Sk, \St, \Se)$ are passed to the shader
as uniforms and determine the displacement amplitudes $A_k$ and
frequencies $f_k$. The quantity $\|\bar{\mathbf{S}}\|^2 = \Sk^2 +
\St^2 + \Se^2$ is computed from these same uniforms. Therefore
$w = e^{-\beta\|\bar{\mathbf{S}}\|^2}$ is available without
additional computation.

\textit{Step 3: Interference from superposition.}
If two tracks' displacement fields $D_A$ and $D_B$ are superposed
(either by rendering both or by passing both phase spectra as
uniforms), the resulting intensity $|D_A + D_B|^2$ contains the
interference term $2\,\mathrm{Re}\bigl[\sum_k A_{A,k} A_{B,k}
e^{i(\Phi_{A,k} - \Phi_{B,k})}\bigr]$, which is proportional to
$V_{AB}$.

\textit{Step 4: Identity.}
All three computations (Boltzmann weight, interference visibility,
partition function contribution) derive from the same set of shader
uniforms and sinusoidal evaluations. They are three projections of a
single spectral computation, unified by the Triple Equivalence
\citep{sachikonye2026categorical}.
\end{proof}

\begin{corollary}[Off-Screen Computation]
\label{cor:offscreen}
The fragment shader need not render to the visible framebuffer. An
off-screen render target computes the same spectral quantities without
visual output. The liquid distortion displayed to the user is a
perceptible side-effect of a computation that would occur identically
in an invisible rendering pass.
\end{corollary}


% ═════════════════════════════════════════════════════════════════════
\section{Purpose-Trained Language Model Interface}
\label{sec:purpose}

The thermodynamic framework provides closed-form solutions for
recommendation, but the user speaks natural language, not statistical
mechanics. A translation layer is required.

\subsection{Knowledge Distillation from Categorical Observations}

The Purpose framework performs six-stage knowledge distillation:

\begin{enumerate}[leftmargin=*,itemsep=2pt]
\item \textbf{Structured Extraction}: Each track observation
  $\mathcal{O}_i$ is converted to structured JSON: S-entropy
  statistics, partition histogram, phase spectrum, dominant partition
  depth.

\item \textbf{Knowledge Mapping}: The collection $\{\mathcal{O}_i\}$
  is aggregated into a knowledge graph: which tracks cluster in
  S-entropy space, what categorical properties define each cluster,
  what interference patterns exist between clusters.

\item \textbf{Stratified Query Generation}: Expert queries are
  generated across four depth levels (basic, intermediate, advanced,
  expert) and five query types (factual, analytical, comparative,
  hypothetical, synthesis), following Bloom's taxonomy
  \citep{bloom1956taxonomy}.

\item \textbf{Multi-Teacher Response Enhancement}: Teacher models
  (GPT-4, Claude) generate responses grounded in the knowledge graph,
  enriched with categorical terminology and S-entropy coordinates.

\item \textbf{Curriculum Organisation}: QA pairs are organised into
  progressive difficulty stages for fine-tuning.

\item \textbf{LoRA Fine-Tuning}: A student model is fine-tuned with
  low-rank adaptation \citep{hu2021lora} on the curriculum, producing
  a domain-specific model that understands this user's categorical
  audio library.
\end{enumerate}

\subsection{Prompt-to-Thermodynamic-Operation Translation}

The trained model translates natural-language prompts into
thermodynamic operations:

\begin{definition}[Prompt Interpretation Function]
\label{def:prompt}
Let $q$ be a natural-language query. The Purpose-trained model
computes:
\begin{equation}
  f_{\mathrm{LLM}}(q, \Lib) = \bigl(\mathbf{S}_{\mathrm{target}},\,
  \Delta T,\, \mathbf{c}_{\mathrm{mask}},\,
  \Lib_{\mathrm{sub}}\bigr)
  \label{eq:prompt-function}
\end{equation}
where $\mathbf{S}_{\mathrm{target}}$ is the target S-entropy region,
$\Delta T$ is a temperature shift (positive for exploration, negative
for refinement), $\mathbf{c}_{\mathrm{mask}}$ is a channel mask
specifying which oscillator classes to prioritise, and
$\Lib_{\mathrm{sub}} \subseteq \Lib$ is a sub-ensemble to condition on
(e.g., ``Tuesday evening'' tracks).
\end{equation}
\end{definition}

\begin{example}
\textit{``Something like the last few tracks but more aggressive.''}

The model identifies: $\Lib_{\mathrm{sub}}$ = last 3 tracks;
$\mathbf{S}_{\mathrm{target}}$ = centroid of sub-ensemble with $\Se$
increased by 0.15; $\Delta T = 0$ (no exploration shift);
$\mathbf{c}_{\mathrm{mask}}$ = weight channels 0, 1, 5 (low
structural, bass-mid coupling, energy envelope).

The search then samples from $P(\mathbf{S}\,|\,\Lib_{\mathrm{sub}})$
with the $\Se$ constraint, translated to Spotify API queries
targeting tracks with matching energy and tempo characteristics.
\end{example}

\subsection{Articulation of Interference Results}

The model also operates in reverse: given an interference result
$V_{AB}$ with channel coherence $\mathbf{c}$, it generates a
natural-language explanation:

\begin{equation}
  g_{\mathrm{LLM}}(V_{AB}, \mathbf{c}, \mathcal{O}_A, \mathcal{O}_B)
  = \text{``These tracks share bass architecture and energy}
  \text{dynamics but diverge in harmonic complexity.''}
  \label{eq:articulation}
\end{equation}

The explanation is grounded in the per-channel decomposition
(Theorem~\ref{thm:interpretability}), not in genre labels or subjective
descriptors.


% ═════════════════════════════════════════════════════════════════════
\section{Ensemble Enrichment Dynamics}
\label{sec:enrichment}

\begin{theorem}[Monotonic Information Gain]
\label{thm:enrichment}
Each new track observation $\mathcal{O}_{M+1}$ added to the library
$\Lib$ strictly increases the Fisher information of the ensemble's
Boltzmann distribution, provided $\mathcal{O}_{M+1}$ is not
categorically identical to an existing observation.
\end{theorem}

\begin{proof}
The Fisher information of the Boltzmann distribution with respect to
the temperature parameter $\beta$ is:
\begin{equation}
  I(\beta) = \mathrm{Var}_P\!\bigl[\|\bar{\mathbf{S}}\|^2\bigr]
  = \langle \|\bar{\mathbf{S}}\|^4 \rangle
  - \langle \|\bar{\mathbf{S}}\|^2 \rangle^2
  \label{eq:fisher}
\end{equation}
Adding a non-degenerate observation increases the variance of the
energy distribution (since the new point either extends the support
or increases the density within the support), strictly increasing
$I(\beta)$ by the strict convexity of variance under non-degenerate
extension.
\end{proof}

\begin{corollary}[Self-Improving Recommendations]
\label{cor:self-improving}
The precision of recommendations improves monotonically with library
size. This improvement arises from thermodynamic enrichment (better
statistical characterisation of the ensemble), not from collaborative
filtering (which requires other users' data) or model retraining
(which requires gradient updates).
\end{corollary}

\begin{remark}
The enrichment is immediate: adding a track observation updates $Z$,
$T_{\Lib}$, $F$, and $S_{\mathrm{taste}}$ in $O(1)$ time (incremental
update of running sums). No batch recomputation is required.
\end{remark}


% ═════════════════════════════════════════════════════════════════════
\section{System Architecture}
\label{sec:architecture}

The complete system comprises five components:

\subsection{Component 1: Categorical Observation Pipeline}

Audio from the Web Audio API analyser is processed at each animation
frame:
\begin{enumerate}[leftmargin=*,itemsep=1pt]
\item FFT yields bass, mid, treble, volume band energies.
\item S-entropy coordinates computed via
  Eqs.~\eqref{eq:sk}--\eqref{eq:se}.
\item Partition coordinates derived from S-entropy (partition depth
  $n$ from $\Se$; harmonic order $\ell$ from $\Sk$; phase index $m$
  from $\St$; chirality $s$ from $\Se$ gradient).
\item Phase spectrum accumulated: $\Phi_k \mathrel{+}= \omega_k \cdot
  \tau_p^{(k)} \cdot \Delta t$.
\item Running statistics (mean, variance) updated incrementally.
\end{enumerate}

\subsection{Component 2: Thermodynamic Ensemble Engine}

Maintains the library $\Lib$ and computes ensemble quantities:
\begin{itemize}[leftmargin=*,itemsep=1pt]
\item Temperature $T_{\Lib}$ (Eq.~\eqref{eq:temperature})
\item Partition function $Z(\beta)$
  (Eq.~\eqref{eq:partition-function})
\item Boltzmann distribution $P(\mathbf{S}\,|\,\Lib)$
  (Eq.~\eqref{eq:boltzmann})
\item Free energy $F$ (Eq.~\eqref{eq:free-energy})
\item Taste entropy $S_{\mathrm{taste}}$
  (Eq.~\eqref{eq:taste-entropy})
\item Interference matrix $V_{ij}$ for all pairs in $\Lib$
\end{itemize}

All quantities are updated incrementally in $O(1)$ per new observation
(except the interference matrix, which requires $O(M)$ per new track).

\subsection{Component 3: Interference Matching Shader}

A WebGL fragment shader evaluates
Eq.~\eqref{eq:shader-displacement} at each pixel, simultaneously
rendering the liquid distortion and computing the spectral quantities.
Phase spectra are passed as uniform arrays. The shader outputs
both the visual frame and (via readback or parallel CPU computation)
the Boltzmann weights and visibility scores.

\subsection{Component 4: Purpose-Trained Language Model}

An LLM (GPT-4 via API, or a LoRA-adapted local model) receives:
\begin{itemize}[leftmargin=*,itemsep=1pt]
\item System prompt: ensemble thermodynamic state ($T_{\Lib}$, $Z$,
  $F$, $S_{\mathrm{taste}}$, cluster summary)
\item User library: per-track categorical observations (JSON)
\item Interference matrix: pairwise visibilities
\item User query: natural language
\end{itemize}

Returns: target S-entropy region, temperature shift, channel mask,
sub-ensemble selection, and Spotify/SoundCloud API search terms.

\subsection{Component 5: Streaming API Integration}

Spotify OAuth PKCE provides track search and 30-second preview
playback. The categorical observation pipeline operates on the preview
audio in real time. Candidate tracks are verified via interference
against the library ensemble before being recommended.


% ═════════════════════════════════════════════════════════════════════
\section{Comparison with Existing Systems}
\label{sec:comparison}

\begin{center}
\begin{tabular}{@{}p{2.3cm}ccc@{}}
\toprule
\textbf{Property} & \textbf{Collab.\ Filter} & \textbf{Embeddings}
  & \textbf{This Work} \\
\midrule
Data source & User behaviour & Audio + text & Audio signal \\
Training & $10^6$+ users & $10^6$+ examples & None \\
Cold start & Fails & Degraded & Immediate \\
Interpretability & None & None & Per-channel \\
Cross-genre & Poor & Moderate & Strong \\
Within-genre & Strong & Strong & Moderate \\
Complexity & $O(Md)$ & $O(Md)$ & $O(MN)$ \\
Improves with & More users & More data & More listening \\
\bottomrule
\end{tabular}
\end{center}

The interference approach does not replace collaborative filtering or
learned embeddings. It provides a \emph{complementary} similarity
signal grounded in the physics of the audio signal, capturing
structural relationships that data-driven methods are structurally
incapable of representing.


% ═════════════════════════════════════════════════════════════════════
\section{Discussion}
\label{sec:discussion}

\subsection{The Physics Framing}

The central observation of this work is that once audio is expressed
as a gas ensemble of oscillators---which the CAT framework proves is
not a metaphor but a mathematical identity---every computational
problem in music recommendation becomes a well-posed problem in
statistical mechanics. Finding similar tracks is computing interference
visibility. Modelling taste is computing a partition function.
Recommending music is sampling from a Boltzmann distribution. Exploring
new genres is heating the ensemble. Refining preference is cooling it.

These are not approximate analogies. They are exact correspondences
between operations on categorical audio state and operations in
statistical mechanics, mediated by the Triple Equivalence: oscillatory
$\equiv$ categorical $\equiv$ partitional. The GPU shader that renders
the visual distortion is the same computation that evaluates the
partition function, because both are spectral evaluations of oscillator
ensembles.

\subsection{Limitations}

\textbf{Representation capacity.} The current implementation uses
$N = 8$ oscillator classes. This is sufficient for demonstrating the
framework but may be insufficient for fine-grained discrimination.
Increasing $N$ to 64 or 256 increases capacity at the cost of $O(N)$
per comparison.

\textbf{Duration normalisation.} Phase accumulation grows with track
duration. Comparing 30-second previews against full tracks requires
normalisation by duration or by phase rates rather than cumulative
phases.

\textbf{Semantic content.} The framework encodes structural
properties of the audio signal. Lyrical content, cultural context,
and subjective associations are outside its scope.

\textbf{Empirical validation.} The mathematical framework is derived
from first principles. Empirical validation through listening studies
correlating interference visibility with human similarity judgments is
required and is the subject of ongoing work.

\subsection{Implications}

If the framework is validated empirically, it implies that a
significant component of musical similarity---the structural
component---is computable from the audio signal alone, without
training data, without user behaviour data, and without cultural
metadata. This component is invisible to existing recommendation
systems because they operate on representations that do not encode
categorical structure.

The practical implication is a recommendation system that works for
the first user on the first track, improves monotonically with
listening, provides interpretable similarity judgments, and discovers
structural connections across genre boundaries that no
collaborative-filtering or embedding-based system can find.


% ═════════════════════════════════════════════════════════════════════
\section{Conclusion}
\label{sec:conclusion}

We have derived a music search and recommendation framework from the
single premise that audio signals are bounded oscillatory systems in
finite phase space. From this premise:

\begin{enumerate}[leftmargin=*,itemsep=2pt]
\item Audio tracks are oscillator ensembles with partition coordinates
  and S-entropy trajectories.
\item The listening library is a thermodynamic gas ensemble with
  temperature, partition function, Boltzmann distribution, free energy,
  and entropy.
\item Similarity between tracks is interference visibility between
  their phase spectra.
\item Recommendation is Boltzmann sampling from the ensemble.
\item The GPU shader rendering visual effects simultaneously computes
  the thermodynamic quantities.
\item A Purpose-trained language model translates between natural
  language and thermodynamic operations.
\item The system improves monotonically with use through ensemble
  enrichment.
\end{enumerate}

No training data is required. No cold-start problem exists. Every
similarity judgment is interpretable per oscillator channel. Cross-genre
structural similarity is captured by construction. The system is
specified with sufficient detail for independent implementation.

\vspace{6pt}
\bibliographystyle{plainnat}
\bibliography{references}

\end{document}
